\chapter{Redução de endomorfismos}\label{ch:b2:reduction}

Para entender um endomorfismo, procure as direções que ele apenas
estica. Este capítulo constrói a maquinaria --- \omterm{def:b2:reduction:eigen}{autovalores},
polinômios característico e minimal, o lema da decomposição
em núcleos --- e suas recompensas: critérios de diagonalização e de
triangularização, Cayley--Hamilton, a \omterm{thm:b2:reduction:dunford}{decomposição de Dunford} e o
cálculo de potências e exponenciais de que o
\cref{ch:b2:diffeq} se alimentará. Em todo o capítulo, $E$ é um
$K$-espaço vetorial de dimensão finita ($K = \R$ ou $\C$) e $u \in
\mathcal{L}(E)$, $n = \dim E$.

\section{Autovalores e autovetores}

\begin{definition}\label{def:b2:reduction:eigen}
$\lambda \in K$ é um \emph{autovalor}\index{autovalor} de $u$
quando $u(x) = \lambda x$ para algum $x \neq 0$ (um
\emph{autovetor}); o \emph{autoespaço} é $E_\lambda(u) =
\ker(u - \lambda\,\mathrm{id})$. O conjunto dos autovalores é o
\emph{espectro}\index{espectro} $\operatorname{Sp}(u)$. Um subespaço $F$ é
\emph{estável} quando $u(F) \subseteq F$; os autoespaços são estáveis, e os
subespaços estáveis permitem endomorfismos induzidos $u|_F$.
\end{definition}

\begin{theorem}[Independência dos autoespaços]\label{thm:b2:reduction:independence}
\omterm{def:b2:reduction:eigen}{Autovetores} associados a \omterm{def:b2:reduction:eigen}{autovalores} dois a dois distintos formam uma
família livre; equivalentemente, a soma dos \omterm{def:b2:reduction:eigen}{autoespaços}
$E_{\lambda_1} + \dots + E_{\lambda_r}$ ($\lambda_i$ distintos) é
direta. Em particular, $u$ tem no máximo $n$ \omterm{def:b2:reduction:eigen}{autovalores}.
\end{theorem}

\begin{proof}
Por indução sobre $r$. Suponha $x_1 + \dots + x_r = 0$ com $x_i \in
E_{\lambda_i}$, sendo o enunciado conhecido para $r - 1$. Aplique $u$ e
subtraia $\lambda_r$ vezes a relação:
\[
\sum_{i=1}^{r-1} (\lambda_i - \lambda_r)\, x_i = 0 ,
\]
logo, por indução, cada $(\lambda_i - \lambda_r)x_i = 0$, isto é, $x_i =
0$ para $i < r$, e então $x_r = 0$. Somas diretas de espaços não nulos num
espaço de dimensão $n$ têm no máximo $n$ parcelas.
\end{proof}

\begin{omfigure}
\begin{tikzpicture}[scale=0.9, >={Stealth}]
  \draw[gray!40] (-2.6,-2.2) grid (2.6,2.2);
  \draw[->] (-2.6,0) -- (2.7,0) node[below, font=\small] {$x$};
  \draw[->] (0,-2.2) -- (0,2.3) node[left, font=\small] {$y$};
  \draw[omDef, very thick, ->] (0,0) -- (0.7,0.7)
    node[above, font=\small] {$v_1$};
  \draw[omDef, thick, ->, dashed] (0,0) -- (2.1,2.1)
    node[above right, font=\small] {$Av_1 = 3v_1$};
  \draw[omThm, very thick, ->] (0,0) -- (0.9,-0.9)
    node[below, font=\small] {$v_2$};
  \draw[omThm, thick, ->, dashed] (0,0) -- (0.9,-0.9);
  \draw[omProp, very thick, ->] (0,0) -- (1,0)
    node[below right, font=\small] {$e_1$};
  \draw[omProp, thick, ->, dashed] (0,0) -- (2,1)
    node[right, font=\small] {$Ae_1$};
\end{tikzpicture}

{\small A matriz $A = \begin{psmallmatrix}2 & 1\\ 1 &
2\end{psmallmatrix}$ agindo no plano: o vetor genérico
$e_1$ é tirado de sua reta, mas as direções próprias $v_1 =
(1,1)$ e $v_2 = (1,-1)$ são apenas esticadas --- por $3$ e por
$1$ (de modo que $Av_2 = v_2$: a imagem tracejada coincide com $v_2$).
A diagonalização é a mudança para a base $(v_1, v_2)$, na qual
$A$ se torna $\operatorname{diag}(3, 1)$.}
\end{omfigure}

\begin{definition}[Polinômio característico]\label{def:b2:reduction:charpoly}
$\chi_u(X) = \det(X\,\mathrm{id} - u)$\index{polinômio
característico} --- calculado em qualquer base como $\det(XI_n - A)$, um polinômio
mônico de grau $n$, invariante por semelhança
(\cref{thm:b2:linalg:detrules}). Suas raízes em $K$ são exatamente os
\omterm{def:b2:reduction:eigen}{autovalores} ($\lambda$ é \omterm{def:b2:reduction:eigen}{autovalor} $\iff u - \lambda\,\mathrm{id}$
não injetivo $\iff \chi_u(\lambda) = 0$), e
\[
\chi_u(X) = X^n - (\operatorname{tr} u)\, X^{n-1} + \dots +
(-1)^n \det u .
\]
A \emph{multiplicidade algébrica}\index{multiplicidade algébrica}
$m_\lambda$ de um \omterm{def:b2:reduction:eigen}{autovalor} é
sua multiplicidade como raiz de $\chi_u$; a \emph{multiplicidade
geométrica}\index{multiplicidade geométrica} é $\dim E_\lambda$, e $1 \leq \dim E_\lambda \leq
m_\lambda$.
\end{definition}

\begin{proof}[Demonstração dos fatos enunciados]
As afirmações sobre os coeficientes: desenvolva $\det(XI - A)$ pela fórmula das
permutações; a permutação identidade contribui com $\prod_i (X - a_{ii})
= X^n - (\sum a_{ii})X^{n-1} + \dots$, e toda outra permutação
fixa no máximo $n - 2$ posições diagonais, contribuindo com grau $\leq
n - 2$: os dois coeficientes de topo são os anunciados; $X = 0$ dá o
termo constante $\det(-A) = (-1)^n\det A$.

Geométrica $\leq$ algébrica: seja $d = \dim E_\lambda$ e complete uma
base de $E_\lambda$ numa base de $E$; a matriz de $u$ é triangular
superior por blocos com bloco superior esquerdo $\lambda I_d$, logo $\chi_u(X) =
(X - \lambda)^d\, \chi_{\text{(bloco inferior)}}(X)$: a multiplicidade
de $\lambda$ vale ao menos $d$.
\end{proof}

\begin{example}[Mesmo $\chi$, geometria diferente]\label{ex:b2:reduction:samechi}
As matrizes
\[
\begin{pmatrix}2 & 0\\ 0 & 2\end{pmatrix}
\qquad\text{e}\qquad
\begin{pmatrix}2 & 1\\ 0 & 2\end{pmatrix}
\]
têm o mesmo \omterm{def:b2:reduction:charpoly}{polinômio característico} $(X - 2)^2$, o mesmo traço, o mesmo
\omterm{def:b2:linalg:det}{determinante}, o mesmo \omterm{def:b2:reduction:eigen}{espectro} --- e no entanto não são semelhantes: a primeira tem
$E_2$ de dimensão $2$ (\omterm{def:b2:reduction:charpoly}{multiplicidade geométrica} $2$), a segunda
de dimensão $1$. O \omterm{def:b2:reduction:charpoly}{polinômio característico} enxerga apenas
\omterm{def:b2:reduction:charpoly}{multiplicidades algébricas}; as dimensões dos \omterm{def:b2:reduction:eigen}{autoespaços} são o
invariante mais fino, e o \omterm{def:b2:reduction:polyu}{polinômio minimal} arbitra ($X - 2$
contra $(X - 2)^2$). Moral para toda discussão de
\omterm{def:b2:reduction:diag}{diagonalizabilidade}: $\chi$ seleciona os candidatos, mas são os núcleos que
depositam os votos.
\end{example}

\begin{definition}[Diagonalizável, triangularizável]\label{def:b2:reduction:diag}
$u$ é \emph{diagonalizável}\index{diagonalizável} quando $E$ tem uma
base de \omterm{def:b2:reduction:eigen}{autovetores} (matricialmente: semelhante a uma matriz diagonal);
\emph{triangularizável}\index{triangularizável} quando sua matriz em alguma
base é triangular superior.
\end{definition}

\begin{theorem}[Critérios de diagonalizabilidade]\label{thm:b2:reduction:diagcrit}
As afirmações seguintes são equivalentes:
\begin{enumerate}
  \item $u$ é \omterm{def:b2:reduction:diag}{diagonalizável};
  \item $E = \bigoplus_{\lambda \in \operatorname{Sp} u} E_\lambda$;
  \item $\chi_u$ se decompõe sobre $K$ e $\dim E_\lambda = m_\lambda$
        para todo \omterm{def:b2:reduction:eigen}{autovalor};
  \item (suficiente, não necessária) $\chi_u$ tem $n$ raízes distintas
        em $K$.
\end{enumerate}
\end{theorem}

\begin{proof}
(1 $\iff$ 2): uma base de \omterm{def:b2:reduction:eigen}{autovetores} se separa em bases dos
$E_\lambda$ e, reciprocamente, concatenar bases das parcelas
diretas dá uma base de $E$
(o~\cref{thm:b2:reduction:independence} torna a soma direta; a igualdade
de dimensões faz dela o espaço todo).

(2 $\iff$ 3): na base diagonal, $\chi_u = \prod (X -
\lambda)^{\dim E_\lambda}$ se decompõe com multiplicidades correspondentes.
Reciprocamente, suponha que $\chi_u$ se decomponha com $\dim E_\lambda =
m_\lambda$ em toda parte; então a soma direta dos \omterm{def:b2:reduction:eigen}{autoespaços}
(direta pelo~\cref{thm:b2:reduction:independence}) tem dimensão
\[
\sum_{\lambda}\dim E_\lambda = \sum_{\lambda} m_\lambda =
\deg\chi_u = n ,
\]
a igualdade do meio porque o grau de um polinômio que se decompõe é a
soma das multiplicidades de suas raízes: a soma é todo o $E$. Note
onde cada hipótese trabalhou: a decomposição preencheu o grau,
a igualdade das multiplicidades preencheu as dimensões.

(4 $\Rightarrow$ 1): $n$ \omterm{def:b2:reduction:eigen}{autovalores} distintos dão $n$ \omterm{def:b2:reduction:eigen}{autovetores}
independentes (\cref{thm:b2:reduction:independence}): uma base.
\end{proof}

\begin{method}[Decidir a diagonalizabilidade]\label{met:b2:reduction:decide}
Na prática, teste nesta ordem --- cada etapa pode encerrar o
serviço. (1) Um polinômio \omterm{def:b2:linalg:annihilator}{anulador} com raízes simples e
que se decompõe se apresenta sozinho ($u^2 = \mathrm{id}$, $u^2 = u$, $u^k =
\mathrm{id}$)? Se sim: \omterm{def:b2:reduction:diag}{diagonalizável}, sem cálculo algum
(\cref{cor:b2:reduction:minpolycrit} adiante). (2) Calcule
$\chi_u$; se ele tiver $n$ raízes distintas em $K$: \omterm{def:b2:reduction:diag}{diagonalizável}
(\cref{thm:b2:reduction:diagcrit}~(4)). (3) Caso contrário, apenas para cada raiz
múltipla $\lambda$, compare $\dim\ker(u -
\lambda\,\mathrm{id})$ com a multiplicidade $m_\lambda$: qualquer
déficit mata a \omterm{def:b2:reduction:diag}{diagonalizabilidade}; a igualdade em toda parte a demonstra.
Nunca calcule \omterm{def:b2:reduction:eigen}{autoespaços} de raízes simples (sua dimensão está
forçada a ser $1$), e nunca triangularize só para decidir.
\end{method}

\begin{example}[A diagonalização posta a trabalhar]\label{ex:b2:reduction:diagwork}
$A = I + J = \begin{psmallmatrix}2 & 1 & 1\\ 1 & 2 & 1\\ 1 & 1 &
2\end{psmallmatrix}$, com $J$ a matriz de uns: de
$\operatorname{Sp}(J) = \{3, 0\}$
(\cref{ex:b2:linalg:onesmatrix}), $\operatorname{Sp}(A) = \{4,
1\}$, com \omterm{def:b2:reduction:eigen}{autoespaços} $\R(1,1,1)$ e o plano $\{x + y + z =
0\}$: dimensões $1 + 2 = 3$, \omterm{def:b2:reduction:diag}{diagonalizável}
(\cref{thm:b2:reduction:diagcrit}~(2)). Potências sem nenhuma
matriz de mudança de base: com $\Pi = J/3$ o projetor sobre
$\R(1,1,1)$,
\[
A = 4\,\Pi + 1\cdot(I - \Pi)
\quad\Longrightarrow\quad
A^k = 4^k\,\Pi + (I - \Pi)
= \frac{4^k - 1}{3}\,J + I .
\]
(Confira $k = 1$: $\frac{4-1}3 J + I = A$.) A lição final:
quando os \omterm{def:b2:reduction:eigen}{autoespaços} estão visíveis, os \emph{projetores} espectrais
calculam potências mais depressa do que $PDP^{-1}$ jamais fará --- e a
fórmula exibe a dinâmica: $A^k$ cresce como $4^k$ ao longo de
$(1,1,1)$ e fica parado no plano ortogonal.
\end{example}

\begin{theorem}[Triangularização]\label{thm:b2:reduction:trigonalization}
$u$ é \omterm{def:b2:reduction:diag}{triangularizável} sobre $K$ se e somente se $\chi_u$ se decompõe sobre
$K$. Em particular, todo endomorfismo de um $\C$-espaço vetorial é
\omterm{def:b2:reduction:diag}{triangularizável}.\index{triangularização}
\end{theorem}

\begin{proof}
($\Rightarrow$) O \omterm{def:b2:reduction:charpoly}{polinômio característico} de uma matriz triangular
é $\prod(X - t_{ii})$: decompõe-se.

($\Leftarrow$) Indução sobre $n$. Como $\chi_u$ se decompõe, ele tem uma
raiz $\lambda$: escolha um \omterm{def:b2:reduction:eigen}{autovetor} $e_1$. Numa base que começa por
$e_1$, a matriz é $\begin{pmatrix} \lambda & \ast\\ 0 &
B\end{pmatrix}$, e $\chi_u = (X - \lambda)\chi_B$: $\chi_B$ também
se decompõe. Pela hipótese de indução aplicada à matriz $(n-1) \times
(n-1)$ $B$, existe uma $Q$ invertível com $Q^{-1}BQ$ triangular
superior; conjugar a matriz inteira por $\begin{pmatrix}1 & 0\\
0 & Q\end{pmatrix}$ a triangulariza.
\end{proof}

\begin{example}[Triangularizando à mão]\label{ex:b2:reduction:trigwork}
$B = \begin{pmatrix}3 & -1\\ 1 & 1\end{pmatrix}$: $\chi_B = X^2 -
4X + 4 = (X - 2)^2$, e $\ker(B - 2I) =
\ker\begin{psmallmatrix}1 & -1\\ 1 & -1\end{psmallmatrix}$ é
a reta \omterm{def:b2:structures:generated}{gerada} por $e_1' = (1, 1)$: um único \omterm{def:b2:reduction:eigen}{autovalor}, um
\omterm{def:b2:reduction:eigen}{autoespaço} de dimensão um --- não \omterm{def:b2:reduction:diag}{diagonalizável}, mas
\omterm{def:b2:reduction:diag}{triangularizável} (\cref{thm:b2:reduction:trigonalization}).
Complete a base com $e_2' = (1, 0)$ e calcule:
\[
u(e_1') = (2, 2) = 2e_1',
\qquad
u(e_2') = (3, 1) = 1\cdot e_1' + 2\, e_2' ,
\]
de modo que, na base $(e_1', e_2')$, a matriz é $T =
\begin{psmallmatrix}2 & 1\\ 0 & 2\end{psmallmatrix}$. A
lição final: a diagonal de $T$ estava forçada (as duas entradas
têm de ser o \omterm{def:b2:reduction:eigen}{autovalor} duplo $2$); só a entrada do canto
dependia da escolha de $e_2'$, e reescalar $e_2'$ pode torná-la
qualquer valor não nulo --- o resistente ``$1$'' é a sombra da
parte nilpotente que Dunford vai isolar.
\end{example}

\section{Polinômios de um endomorfismo}

\begin{definition}\label{def:b2:reduction:polyu}
Para $P = \sum a_k X^k \in K[X]$, ponha $P(u) = \sum a_k u^k \in
\mathcal{L}(E)$. A aplicação $P \mapsto P(u)$ é um morfismo de \omterm{def:b2:structures:algebra}{álgebras}
$K[X] \to \mathcal{L}(E)$ (\cref{def:b2:structures:algebra});
seu núcleo $\{P : P(u) = 0\}$ é um \omterm{def:b2:structures:ideal}{ideal} de $K[X]$, não nulo (a
família $(\mathrm{id}, u, \dots, u^{n^2})$ é ligada no
$n^2$ de dimensão $\mathcal{L}(E)$), logo \omterm{def:b2:structures:generated}{gerado} por um único
polinômio mônico $\mu_u$: o \emph{polinômio
minimal}\index{polinômio minimal}
(\cref{thm:b2:structures:principal}).
\end{definition}

\begin{proposition}\label{prop:b2:reduction:minpoly}
\begin{enumerate}
  \item $P(u) = 0 \iff \mu_u \mid P$; os \omterm{def:b2:reduction:eigen}{autovalores} de $u$ são raízes
        de todo polinômio \omterm{def:b2:linalg:annihilator}{anulador}, e as raízes de $\mu_u$
        são \emph{exatamente} os \omterm{def:b2:reduction:eigen}{autovalores}.
  \item Se $F$ é estável, $\mu_{u|_F} \mid \mu_u$.
\end{enumerate}
\end{proposition}

\begin{proof}
(1) A divisibilidade é a definição de gerador. Se $u(x) =
\lambda x$, $x \neq 0$, então $0 = P(u)(x) = P(\lambda)x$, logo
$P(\lambda) = 0$: os \omterm{def:b2:reduction:eigen}{autovalores} são raízes dos \omterm{def:b2:linalg:annihilator}{anuladores}, em
particular de $\mu_u$. Reciprocamente, se $\lambda$ é raiz, $\mu_u =
(X - \lambda)Q$ com $Q(u) \neq 0$ (o grau de $\mu_u$ é mínimo): escolha
$y$ com $Q(u)(y) \neq 0$; então $(u - \lambda)(Q(u)(y)) = \mu_u(u)(y)
= 0$ exibe o \omterm{def:b2:reduction:eigen}{autovetor} $Q(u)(y)$.

(2) $\mu_u(u|_F) = \mu_u(u)|_F = 0$, e aplique (1) a $u|_F$.
\end{proof}

\begin{example}[Polinômios minimais achados à mão]\label{ex:b2:reduction:minpolyhand}
O \omterm{def:b2:reduction:polyu}{polinômio minimal} se calcula testando graus
sucessivos. Para a matriz de uns $J \in \mathcal{M}_3(\R)$:
$J \neq \lambda I$ (o grau $1$ está fora), e $J^2 = 3J$, logo
\[
\mu_J = X^2 - 3X = X(X - 3) :
\]
grau $2$, decompõe-se, raízes simples --- $J$ é \omterm{def:b2:reduction:diag}{diagonalizável} com
\omterm{def:b2:reduction:eigen}{espectro} $\{0, 3\}$ (\cref{cor:b2:reduction:minpolycrit} adiante),
confirmando \cref{ex:b2:linalg:onesmatrix} sem um único
\omterm{def:b2:linalg:det}{determinante}. Para a matriz de troca $A$ de
\cref{ex:b2:reduction:projectorswork}: $A \neq \pm I$ e $A^2
= I$ dão $\mu_A = X^2 - 1$. Nos dois casos o padrão é o
mesmo: adivinhe uma identidade de grau baixo a partir da estrutura
(o posto um força $J^2 = (\operatorname{tr}J)\,J$; uma involução
força $A^2 = I$) e depois verifique que nenhum divisor próprio
anula. Os \omterm{def:b2:reduction:polyu}{polinômios minimais} são em geral \emph{achados}, e não
calculados a partir de $\chi$.
\end{example}

\begin{theorem}[Lema da decomposição em núcleos]\label{thm:b2:reduction:kernels}
Se $P = P_1 P_2 \cdots P_r$ com os $P_i$ dois a dois coprimos,
então\index{lema da decomposição em núcleos}
\[
\ker P(u) = \ker P_1(u) \oplus \dots \oplus \ker P_r(u),
\]
e as projeções sobre as parcelas são polinômios em $u$.
\end{theorem}

\begin{proof}
Basta tratar o caso $r = 2$ e fazer indução. Bézout em $K[X]$
(\cref{thm:b2:structures:principal}): $U P_1 + V P_2 = 1$, de modo que, para
todo $x$,
\[
x = \underbrace{U(u)P_1(u)(x)}_{=:\,x_2}
+ \underbrace{V(u)P_2(u)(x)}_{=:\,x_1}.
\]
Se $x \in \ker P(u)$: $P_2(u)(x_2) = U(u)\,P(u)(x) = 0$ (polinômios
em $u$ comutam), logo $x_2 \in \ker P_2(u)$ e, simetricamente, $x_1
\in \ker P_1(u)$: a soma preenche $\ker P(u)$; as duas parcelas ficam
dentro de $\ker P(u)$ ($P_i \mid P$). Diretividade: $x \in \ker P_1(u)
\cap \ker P_2(u)$ dá $x = U(u)P_1(u)x + V(u)P_2(u)x = 0$. As
fórmulas para $x_1, x_2$ exibem as projeções como $V(u)P_2(u)$ e
$U(u)P_1(u)$.
\end{proof}

\begin{example}[O lema dos núcleos com projetores explícitos]\label{ex:b2:reduction:projectorswork}
Seja $A = \begin{psmallmatrix}0 & 1 & 0\\ 1 & 0 & 0\\ 0 & 0 &
1\end{psmallmatrix}$ (troca as duas primeiras coordenadas). Então $A^2
= I$: o polinômio $X^2 - 1 = (X - 1)(X + 1)$ anula $A$,
seus fatores são coprimos e Bézout é explícito:
\[
\frac{1}{2}(X + 1) - \frac12(X - 1) = 1 .
\]
Seguindo a demonstração do~\cref{thm:b2:reduction:kernels}, as
projeções sobre $\ker(A - I)$ e $\ker(A + I)$ são os
polinômios em $A$
\[
\pi_+ = \frac{A + I}{2} = \frac12\begin{pmatrix}
1 & 1 & 0\\ 1 & 1 & 0\\ 0 & 0 & 2\end{pmatrix},
\qquad
\pi_- = \frac{I - A}{2} = \frac12\begin{pmatrix}
1 & -1 & 0\\ -1 & 1 & 0\\ 0 & 0 & 0\end{pmatrix}.
\]
Verificação: $\pi_+ + \pi_- = I$, $\pi_+\pi_- = 0$, $\pi_\pm^2 =
\pi_\pm$, e as imagens são o plano $\{x = y\}$ (vetores
simétricos, \omterm{def:b2:reduction:eigen}{autovalor} $1$) e a reta $\R(1, -1, 0)$
(antissimétricos, \omterm{def:b2:reduction:eigen}{autovalor} $-1$). O lema dos núcleos não é um
enunciado de existência: os coeficientes de Bézout \emph{são} as
fórmulas dos projetores.
\end{example}

\begin{example}[Os projetores também calculam a exponencial]\label{ex:b2:reduction:expprojectors}
A mesma matriz de troca, um dividendo adiante. Como $A = \pi_+ -
\pi_-$ com projetores ortogonais no sentido algébrico
($\pi_+\pi_- = 0$), toda potência obedece a $A^k = \pi_+ +
(-1)^k\pi_-$, e a série exponencial se reagrupa por projetor:
\[
\eu^{tA} = \sum_k \frac{t^k}{k!}\bigl(\pi_+ +
(-1)^k\pi_-\bigr)
= \eu^{t}\,\pi_+ + \eu^{-t}\,\pi_- =
\begin{pmatrix}
\cosh t & \sinh t & 0\\
\sinh t & \cosh t & 0\\
0 & 0 & \eu^{t}
\end{pmatrix}.
\]
(Confira $t = 0$: a identidade; derivada em $0$: $A$.) A
decomposição espectral converte uma série de matrizes em duas séries
escalares --- exatamente o mecanismo que o~\cref{ch:b2:diffeq} vai
acionar em todo sistema \omterm{def:b2:reduction:diag}{diagonalizável}, e a razão pela qual as funções
hiperbólicas governam acoplamentos simétricos.
\end{example}

\begin{corollary}[Diagonalizabilidade pelo polinômio minimal]\label{cor:b2:reduction:minpolycrit}
$u$ é \omterm{def:b2:reduction:diag}{diagonalizável} $\iff$ $\mu_u$ se decompõe sobre $K$ com
raízes \emph{simples} $\iff$ algum polinômio \omterm{def:b2:linalg:annihilator}{anulador} de $u$
se decompõe com raízes simples.
\end{corollary}

\begin{proof}
Se $P(u) = 0$ com $P = \prod_{i}(X - \lambda_i)$ ($\lambda_i$
distintos), o lema dá $E = \ker P(u) = \bigoplus_i \ker(u -
\lambda_i)$: uma soma direta de \omterm{def:b2:reduction:eigen}{autoespaços}, logo $u$ é \omterm{def:b2:reduction:diag}{diagonalizável}
(\cref{thm:b2:reduction:diagcrit}). Reciprocamente, um $u$ \omterm{def:b2:reduction:diag}{diagonalizável}
é anulado por $\prod_{\lambda \in \operatorname{Sp}u}(X -
\lambda)$ (que anula cada \omterm{def:b2:reduction:eigen}{autoespaço}), o qual se decompõe com raízes
simples; e $\mu_u$ o divide tendo as mesmas raízes
(\cref{prop:b2:reduction:minpoly}): $\mu_u$ é exatamente esse produto.
\end{proof}

\begin{example}\label{ex:b2:reduction:projections}
As projeções satisfazem $p^2 = p$: anuladas por $X(X-1)$, que se decompõe com raízes
simples --- \omterm{def:b2:reduction:diag}{diagonalizáveis} com \omterm{def:b2:reduction:eigen}{espectro} $\subseteq \{0, 1\}$, e $E =
\ker p \oplus \ker(p - \mathrm{id})$: a análise geométrica do primeiro ano,
redemonstrada em uma linha. As simetrias ($s^2 = \mathrm{id}$,
\omterm{def:b2:linalg:annihilator}{anulador} $X^2 - 1$): \omterm{def:b2:reduction:diag}{diagonalizáveis} quando $\operatorname{char} K
\neq 2$, \omterm{def:b2:reduction:eigen}{espectro} $\subseteq \{\pm 1\}$. Um endomorfismo com $u^3 =
u^2$ e $u^2 \neq u$: anulado por $X^2(X - 1)$, \emph{não}
necessariamente \omterm{def:b2:reduction:diag}{diagonalizável} --- o critério o detecta (a raiz dupla
$0$ precisa ser testada: \omterm{def:b2:reduction:diag}{diagonalizável} se e somente se além disso $\ker u^2 = \ker
u$).
\end{example}

\begin{example}[O corpo decide: uma rotação em $\R^3$]\label{ex:b2:reduction:rotation3d}
Seja $R$ o quarto de volta em torno do eixo $z$:
\[
R = \begin{pmatrix}
0 & -1 & 0\\
1 & 0 & 0\\
0 & 0 & 1
\end{pmatrix},
\qquad
\chi_R = (X - 1)(X^2 + 1).
\]
Sobre $\R$: o único \omterm{def:b2:reduction:eigen}{autovalor} é $1$, com \omterm{def:b2:reduction:eigen}{autoespaço} igual ao eixo
$\R e_3$ --- uma reta de vetores fixos e nenhuma outra
redução: $R$ não é \omterm{def:b2:reduction:diag}{diagonalizável} nem \omterm{def:b2:reduction:diag}{triangularizável} em
$\mathcal{M}_3(\R)$ ($\chi_R$ não se decompõe). Sobre $\C$: três
\omterm{def:b2:reduction:eigen}{autovalores} distintos $1, \iu, -\iu$, de modo que $R$ é \omterm{def:b2:reduction:diag}{diagonalizável},
com \omterm{def:b2:reduction:eigen}{autovetores} $e_3$ e $e_1 \mp \iu e_2$. A geometria era
audível na álgebra: as rotações do plano não têm direções invariantes
reais, e os \omterm{def:b2:reduction:eigen}{autovalores} complexos $\pm\iu$ de
módulo $1$ guardam o ângulo ($\pm\frac\pi2$) que a matriz real
só consegue exprimir misturando coordenadas.
\end{example}

\begin{example}[Minimal contra característico]\label{ex:b2:reduction:minvschar}
Para $D = \operatorname{diag}(2, 2, 3)$: $\chi_D = (X - 2)^2(X -
3)$ mas $\mu_D = (X - 2)(X - 3)$, pois $(D - 2I)(D - 3I) = 0$
(verifique na base canônica), ao passo que nenhum dos fatores sozinho anula
$D$. Para o bloco de deslocamento $N = \begin{psmallmatrix}0 & 1\\ 0 &
0\end{psmallmatrix} \oplus (3)$, isto é, $N' =
\begin{psmallmatrix}0 & 1 & 0\\ 0 & 0 & 0\\ 0 & 0 &
3\end{psmallmatrix}$: $\chi_{N'} = X^2(X - 3)$ \emph{e}
$\mu_{N'} = X^2(X - 3)$ --- a raiz dupla é genuinamente necessária,
pois $N'$ não é \omterm{def:b2:reduction:diag}{diagonalizável} do lado $\ker$ ($N'e_2 =
e_1 \neq 0$). Regra prática: $\mu$ e $\chi$ têm as mesmas raízes
(\cref{prop:b2:reduction:minpoly}); a multiplicidade em $\mu$
mede o tamanho do maior bloco nilpotente, e a de
$\chi$ a dimensão total do subespaço característico.
\end{example}

\begin{theorem}[Cayley--Hamilton]\label{thm:b2:reduction:cayleyhamilton}
$\chi_u(u) = 0$; por consequência, $\mu_u \mid \chi_u$ e $\deg \mu_u
\leq n$.\index{teorema de Cayley--Hamilton}
\end{theorem}

\begin{proof}
Fixe $x \neq 0$ e seja $d$ maximal com $(x, u(x), \dots,
u^{d-1}(x))$ livre; escreva
\[
u^d(x) = -a_0 x - a_1 u(x) - \dots - a_{d-1}u^{d-1}(x),
\]
e ponha $P_x = X^d + a_{d-1}X^{d-1} + \dots + a_0$, de modo que $P_x(u)(x) =
0$. Complete a família livre numa base de $E$: nela, $u$ tem
a forma por blocos $\begin{pmatrix} C & \ast\\ 0 & D\end{pmatrix}$, em que $C$
é a matriz companheira de $P_x$, cujo \omterm{def:b2:reduction:charpoly}{polinômio característico} é
$P_x$ (desenvolva $\det(XI - C)$ ao longo da primeira coluna, por indução sobre
$d$). Logo $\chi_u = P_x \cdot \chi_D$, e
\[
\chi_u(u)(x) = \chi_D(u)\bigl(P_x(u)(x)\bigr) = 0 .
\]
O argumento vale para todo $x$: $\chi_u(u) = 0$.
\end{proof}

\begin{example}[Cayley--Hamilton em ação]\label{ex:b2:reduction:chwork}
$A = \begin{pmatrix} 1 & 2\\ 3 & 4\end{pmatrix}$: $\chi_A = X^2 -
5X - 2$, logo $A^2 = 5A + 2I$. Toda potência de $A$ desaba para uma
combinação de $I$ e $A$:
\[
A^4 = (5A + 2I)^2 = 25A^2 + 20A + 4I = 145A + 54I =
\begin{pmatrix} 199 & 290\\ 435 & 634\end{pmatrix},
\]
e a inversa vem de brinde: $A(A - 5I) = 2I$ dá
\[
A^{-1} = \tfrac12(A - 5I)
= \begin{pmatrix} -2 & 1\\ 3/2 & -1/2\end{pmatrix}.
\]
A lição final: Cayley--Hamilton comprime toda a \omterm{def:b2:structures:algebra}{álgebra}
$K[A]$ em $\operatorname{Vect}(I, A, \dots, A^{n-1})$ ---
$\dim K[A] = \deg\mu_A \leq n$, por maiores que sejam as potências de que você precise.
\end{example}

\begin{remark}[Armadilhas comuns]\label{rem:b2:reduction:pitfalls}
(i) \omterm{def:b2:reduction:eigen}{Autovalores} não se somam: $\operatorname{Sp}(A + B)$ não é
$\operatorname{Sp}A + \operatorname{Sp}B$, e uma soma de matrizes
\omterm{def:b2:reduction:diag}{diagonalizáveis} não precisa ser \omterm{def:b2:reduction:diag}{diagonalizável} ---
$\begin{psmallmatrix}1 & 1\\ 0 & 0\end{psmallmatrix} +
\begin{psmallmatrix}0 & 0\\ 0 & 1\end{psmallmatrix} =
\begin{psmallmatrix}1 & 1\\ 0 & 1\end{psmallmatrix}$ é soma de
duas matrizes \omterm{def:b2:reduction:diag}{diagonalizáveis} (cada uma com \omterm{def:b2:reduction:eigen}{autovalores} distintos) e
não é \omterm{def:b2:reduction:diag}{diagonalizável}; só as famílias \emph{que comutam} se comportam bem
(\cref{exo:b2:reduction:9}). (ii) ``$\chi_u$ se decompõe'' é uma
hipótese sobre o \emph{corpo}: uma rotação plana tem $\chi =
X^2 - 2\cos\theta\,X + 1$, que se decompõe sobre $\C$, não sobre $\R$ ---
\omterm{def:b2:reduction:diag}{diagonalizável} em $\mathcal{M}_2(\C)$, não \omterm{def:b2:reduction:diag}{triangularizável} em
$\mathcal{M}_2(\R)$. (iii) A desigualdade vai da geométrica $\leq$
à algébrica, nunca ao contrário; testar apenas $\dim E_\lambda \geq
1$ nada prova sobre a \omterm{def:b2:reduction:diag}{diagonalizabilidade}. (iv) $\mu_u$ não é
$\chi_u$: a igualdade vale exatamente quando cada \omterm{def:b2:reduction:eigen}{autovalor} tem uma
única cadeia de blocos (por exemplo, as matrizes companheiras, no
problema de fim de semana deste capítulo); usar $\chi$ onde $\mu$ é necessário infla
todo cálculo de potências. (v) As partes $d$ e $\nu$ de Dunford são
polinômios em $u$ --- uma decomposição $u = d' + \nu'$ com as
propriedades certas mas com $d'\nu' \neq \nu'd'$ \emph{não} é a de Dunford
e nunca é única.
\end{remark}

\begin{remark}[Onde este capítulo é usado]
A redução é o cavalo de batalha do resto do livro: potências e
exponenciais de matrizes movem os sistemas diferenciais lineares do
\cref{ch:b2:diffeq}; o teorema espectral do~\cref{ch:b2:quadratic}
é a diagonalização tornada ortogonal; as funções geradoras
(\cref{ch:b2:genfun}) redemonstram analiticamente as assintóticas de recorrência do
problema de fim de semana deste capítulo. No volume do terceiro ano de graduação, o
mesmo programa roda em dimensão infinita: a teoria espectral dos
operadores compactos autoadjuntos, na qual sequências de \omterm{def:b2:reduction:eigen}{autovalores} substituem
os \omterm{def:b2:reduction:eigen}{espectros} finitos, e a teoria de Perron--Frobenius das matrizes
positivas, que explica \emph{por que} os \omterm{def:b2:reduction:eigen}{autovalores} dominantes de
problemas de contagem são positivos e simples.
\end{remark}

\section{Nilpotentes e a decomposição de Dunford}

\begin{proposition}[Endomorfismos nilpotentes]\label{prop:b2:reduction:nilpotent}
Para $u$ com $\chi_u$ que se decompõe, as afirmações seguintes são equivalentes: $u^n =
0$; $u^k = 0$ para algum $k$; $\operatorname{Sp}(u) = \{0\}$; $\chi_u
= X^n$; $u$ é \omterm{def:b2:reduction:diag}{triangularizável} com diagonal nula. Um \omterm{prop:b2:reduction:nilpotent}{endomorfismo
nilpotente} tem $\mu_u = X^{\text{(índice de nilpotência)}}$, e o índice satisfaz
$\leq n$.\index{endomorfismo nilpotente}
\end{proposition}

\begin{proof}
$u^k = 0$ faz de todo \omterm{def:b2:reduction:eigen}{autovalor} uma raiz de $X^k$: \omterm{def:b2:reduction:eigen}{espectro} $\{0\}$
(não vazio quando $\chi$ se decompõe --- sobre $\C$ sempre). Então $\chi_u =
X^n$ (todas as raízes nulas) e Cayley--Hamilton dá $u^n = 0$;
a triangularização (\cref{thm:b2:reduction:trigonalization}) põe
zeros na diagonal (a diagonal carrega os \omterm{def:b2:reduction:eigen}{autovalores}).
Reciprocamente, seja $A$ estritamente triangular superior: $a_{ij} = 0$
para $j \leq i$. Mostramos por indução que
\[
(A^k)_{ij} = 0 \qquad \text{sempre que } j \leq i + k - 1,
\]
isto é, cada potência empurra a região nula uma diagonal acima. Para
$k = 1$ essa é a hipótese. Para o passo,
\[
(A^{k+1})_{ij} = \sum_{\ell} (A^k)_{i\ell}\,a_{\ell j} ,
\]
e cada termo se anula: ou $\ell \leq i + k - 1$ (o primeiro
fator é $0$ por indução), ou $\ell \geq i + k$, e nesse caso
$j \leq i + k \leq \ell$ mata o segundo fator. Em $k = n$ a
condição $j \leq i + n - 1$ vale para todos $i, j \leq n$: $A^n =
0$. O \omterm{def:b2:reduction:polyu}{polinômio minimal} divide $X^n$ e a anulação
define o índice.
\end{proof}

\begin{theorem}[Decomposição de Dunford]\label{thm:b2:reduction:dunford}
Suponha que $\chi_u$ se decomponha sobre $K$ (automático para $K = \C$). Então
existe um \emph{único} par $(d, \nu)$ com
\[
u = d + \nu, \qquad d \text{ diagonalizável}, \quad \nu
\text{ nilpotente}, \quad d\nu = \nu d ,
\]
e além disso $d$ e $\nu$ são polinômios em
$u$.\index{decomposição de Dunford}
\end{theorem}

\begin{proof}
\emph{Existência.} Escreva $\chi_u = \prod_{i=1}^{r} (X -
\lambda_i)^{m_i}$ ($\lambda_i$ distintos) e ponha $N_i = \ker(u -
\lambda_i)^{m_i}$, os \emph{subespaços característicos}. Por
Cayley--Hamilton e pelo lema dos núcleos
(\cref{thm:b2:reduction:kernels}),
\[
E = N_1 \oplus \dots \oplus N_r ,
\]
com projeções $\pi_i$ polinomiais em $u$; cada $N_i$ é estável
(polinômios em $u$ comutam com $u$). Defina $d = \sum_i \lambda_i
\pi_i$: um polinômio em $u$, \omterm{def:b2:reduction:diag}{diagonalizável} (ele age como $\lambda_i$
em $N_i$, logo $E$ se decompõe em seus \omterm{def:b2:reduction:eigen}{autoespaços}). Então $\nu = u -
d$ é um polinômio em $u$ (logo comuta com $d$) e, em cada
$N_i$, age como $u - \lambda_i$, com $(u - \lambda_i)^{m_i} = 0$
aí: $\nu^{\max m_i} = 0$ em cada parcela, logo $\nu$ é nilpotente.

\emph{Unicidade.} Seja $u = d' + \nu'$ outro par desse tipo. Como
$d'$ e $\nu'$ comutam entre si, eles comutam com $u = d' +
\nu'$, logo com todo polinômio em $u$ --- em particular com $d$
e $\nu$. Então $d - d'$ é \omterm{def:b2:reduction:diag}{diagonalizável} (duas aplicações
\omterm{def:b2:reduction:diag}{diagonalizáveis} que comutam são simultaneamente \omterm{def:b2:reduction:diag}{diagonalizáveis}:
\cref{exo:b2:reduction:9}) e é igual a $\nu' - \nu$, que é
nilpotente: se $\nu^k = 0$ e $\nu'^{k'} = 0$, a comutatividade
autoriza o desenvolvimento binomial
\[
(\nu' - \nu)^{k + k' - 1}
= \sum_{j=0}^{k+k'-1}\binom{k + k' - 1}{j}
\,\nu'^{\,j}\,(-\nu)^{k + k' - 1 - j} ,
\]
no qual todo termo morre: ou $j \geq k'$ (primeiro fator
nulo), ou $k + k' - 1 - j \geq k$ (segundo fator nulo), e uma
das duas condições sempre vale. Um nilpotente \omterm{def:b2:reduction:diag}{diagonalizável} é nulo (seu
\omterm{def:b2:reduction:eigen}{espectro} é $\{0\}$ e ele é diagonal em alguma base): $d = d'$,
$\nu = \nu'$.
\end{proof}

\begin{example}[Potências e exponenciais]\label{ex:b2:reduction:powers}
$A = \begin{pmatrix} 3 & 1\\ -1 & 1\end{pmatrix}$: $\chi_A = X^2 -
4X + 4 = (X-2)^2$, um único \omterm{def:b2:reduction:eigen}{autovalor} $2$, \omterm{def:b2:reduction:eigen}{autoespaço} de dimensão
$1$: não \omterm{def:b2:reduction:diag}{diagonalizável}. Dunford: $D = 2I$, $N = A - 2I =
\begin{pmatrix} 1 & 1\\ -1 & -1\end{pmatrix}$, $N^2 = 0$. Então
\[
A^k = (2I + N)^k = 2^k I + k\,2^{k-1} N ,
\qquad
\eu^{tA} = \eu^{2t}(I + tN),
\]
pelo binômio de Newton para elementos que comutam, resp.\ pela série exponencial
(\cref{ch:b2:diffeq}) separada em parcelas que comutam. A redução transforma
dinâmica matricial em dinâmica escalar.
\end{example}

\begin{remark}[Perspectivas dentro deste volume]\label{rem:b2:reduction:perspectives}
A redução é um entroncamento; eis os quatro ramais a observar. No
\cref{ch:b2:nvs}, normas adaptadas transformam ``todos os \omterm{def:b2:reduction:eigen}{autovalores} de
módulo $< 1$'' em ``alguma norma de operador $< 1$'', fazendo os
\omterm{def:b2:reduction:eigen}{espectros} governarem a convergência de potências e séries. No
\cref{ch:b2:diffeq}, a receita do
\cref{ex:b2:reduction:powers} torna-se a solução geral de
$X' = AX$: Dunford separa $\eu^{tA}$ em blocos do tipo
polinômio vezes exponencial, e a estabilidade se lê nas
partes reais dos \omterm{def:b2:reduction:eigen}{autovalores}. No~\cref{ch:b2:quadratic}, um
produto escalar impõe o que a mera álgebra linear não consegue: as matrizes
simétricas tornam-se \emph{ortogonalmente} \omterm{def:b2:reduction:diag}{diagonalizáveis}, sem
nenhuma parte nilpotente. E no~\cref{ch:b2:genfun}, as
assintóticas de \omterm{def:b2:reduction:eigen}{autovalor} dominante do problema de fim de semana
deste capítulo reaparecem analiticamente, como a menor singularidade de uma
função geradora --- duas línguas para uma mesma taxa de crescimento.
\end{remark}

\section{Exercícios}

\begin{exercise}[$\star$]\label{exo:b2:reduction:1}
Diagonalize (\omterm{def:b2:reduction:eigen}{autovalores}, bases dos \omterm{def:b2:reduction:eigen}{autoespaços}, $P$ invertível):
\[
A = \begin{pmatrix} 1 & 2\\ 2 & 1 \end{pmatrix},
\qquad
B = \begin{pmatrix} 0 & 1 & 1\\ 1 & 0 & 1\\ 1 & 1 & 0
\end{pmatrix}.
\]
\end{exercise}

\begin{exercise}[$\star$]\label{exo:b2:reduction:2}
Mostre que $C = \begin{pmatrix} 1 & 1\\ 0 & 1\end{pmatrix}$ não é
\omterm{def:b2:reduction:diag}{diagonalizável}, de duas maneiras: pelos \omterm{def:b2:reduction:eigen}{autoespaços} e pelo \omterm{def:b2:reduction:polyu}{polinômio
minimal}.
\end{exercise}

\begin{exercise}[$\star$]\label{exo:b2:reduction:3}
Seja $u$ tal que $u^2 - 5u + 6\,\mathrm{id} = 0$. Prove que $u$ é
\omterm{def:b2:reduction:diag}{diagonalizável}, determine os \omterm{def:b2:reduction:eigen}{espectros} possíveis e calcule $u^k$
como combinação de $\mathrm{id}$ e $u$.
\end{exercise}

\begin{exercise}[$\star\star$]\label{exo:b2:reduction:4}
Seja $u$ \omterm{def:b2:reduction:diag}{diagonalizável} e $F$ um subespaço estável. Prove que
$u|_F$ é \omterm{def:b2:reduction:diag}{diagonalizável} \emph{(restrinja um polinômio \omterm{def:b2:linalg:annihilator}{anulador}
com raízes simples e que se decompõe)}.
\end{exercise}

\begin{exercise}[$\star\star$]\label{exo:b2:reduction:5}
(Fibonacci) Seja $A = \begin{pmatrix} 1 & 1\\ 1 & 0\end{pmatrix}$.
Diagonalize $A$ sobre $\R$ e deduza a fórmula de Binet para a
sequência de Fibonacci ($F_0 = 0$, $F_1 = 1$, $F_{n+1} = F_n +
F_{n-1}$):
\[
F_n = \frac{\varphi^n - \psi^n}{\sqrt 5},
\qquad \varphi = \frac{1 + \sqrt5}{2},\ \psi = \frac{1 -
\sqrt5}{2}.
\]
\end{exercise}

\begin{exercise}[$\star\star$]\label{exo:b2:reduction:6}
Sejam $u \in \mathcal{L}(E)$ com $u^2$ \omterm{def:b2:reduction:diag}{diagonalizável} e $u$
invertível ($K = \C$). Prove que $u$ é \omterm{def:b2:reduction:diag}{diagonalizável}. Dê um
contraexemplo quando $u$ não é invertível.
\end{exercise}

\begin{exercise}[$\star\star$]\label{exo:b2:reduction:7}
Calcule a \omterm{thm:b2:reduction:dunford}{decomposição de Dunford}, $A^k$ e $\eu^{tA}$ para
\[
A = \begin{pmatrix} 2 & 1 & 0\\ 0 & 2 & 1\\ 0 & 0 & 2
\end{pmatrix}.
\]
\end{exercise}

\begin{exercise}[$\star\star$]\label{exo:b2:reduction:8}
Seja $A \in \mathcal{M}_n(\C)$ com $A^k = I$ para algum $k \geq 1$.
Prove que $A$ é \omterm{def:b2:reduction:diag}{diagonalizável} e que seus \omterm{def:b2:reduction:eigen}{autovalores} são raízes
$k$-ésimas da unidade. Deduza que uma matriz complexa invertível de \omterm{def:b2:structures:generated}{ordem} finita
semelhante a uma matriz triangular com diagonal de uns é a identidade.
\end{exercise}

\begin{exercise}[$\star\star\star$]\label{exo:b2:reduction:9}
(Diagonalização simultânea) Sejam $u, v$ \omterm{def:b2:reduction:diag}{diagonalizáveis} e
que comutam. Prove que eles são \emph{simultaneamente}
\omterm{def:b2:reduction:diag}{diagonalizáveis}: alguma base diagonaliza os dois. \emph{(Cada \omterm{def:b2:reduction:eigen}{autoespaço}
de $u$ é estável por $v$; diagonalize aí as restrições de $v$,
usando \cref{exo:b2:reduction:4}.)}
\end{exercise}

\begin{exercise}[$\star\star\star$]\label{exo:b2:reduction:10}
Seja $u \in \mathcal{L}(\C^n)$. Prove que $u$ é \omterm{def:b2:reduction:diag}{diagonalizável} se
e somente se todo subespaço $u$-estável tem um subespaço suplementar
$u$-estável. \emph{(Para $\Leftarrow$: aplique a propriedade
a $F = \sum_\lambda E_\lambda(u)$, a soma de todos os \omterm{def:b2:reduction:eigen}{autoespaços}; se um
suplementar estável $G$ fosse não nulo, triangularizar $u|_G$ produziria
um \omterm{def:b2:reduction:eigen}{autovetor} de $u$ dentro de $G$ --- contradizendo $G \cap
F = \{0\}$.)}
\end{exercise}

\begin{exercise}[$\star\star\star$]\label{exo:b2:reduction:11}
(Raio espectral à la Gelfand, $2\times2$ aperitivo da análise que
vem) Seja $A \in \mathcal{M}_2(\C)$ com os dois \omterm{def:b2:reduction:eigen}{autovalores} de
módulo $< 1$. Prove que $A^k \to 0$ entrada a entrada quando $k \to \infty$.
\emph{(Triangularize: $A = P(T)P^{-1}$ com $T$ triangular superior;
calcule $T^k$ explicitamente --- distinga \omterm{def:b2:reduction:eigen}{autovalores} iguais e
distintos --- e majore.)}
\end{exercise}

\begin{exercise}[$\star\star$]\label{exo:b2:reduction:12}
Seja $u \in \mathcal{L}(\C^n)$ com $\operatorname{rk} u = 1$ ($n
\geq 2$). Mostre que $\chi_u = X^{n-1}(X - \operatorname{tr} u)$
e que $u$ é \omterm{def:b2:reduction:diag}{diagonalizável} se e somente se $\operatorname{tr} u
\neq 0$. \emph{(Lembre do~\cref{exo:b2:linalg:5} que $u^2 =
(\operatorname{tr} u)\,u$.)}
\end{exercise}

\section{Problema: recorrências lineares e matrizes companheiras}

Uma recorrência linear $u_{n+k} = a_{k-1}u_{n+k-1} + \dots + a_0 u_n$
é uma potência de matriz disfarçada, e a redução a converte em
fórmulas fechadas, taxas de crescimento e estimativas de erro. Este problema
de fim de semana desenvolve o dicionário --- matrizes companheiras de um
lado, o operador de deslocamento no espaço das sequências do outro
--- demonstra o \emph{teorema fundamental das recorrências lineares}
(a solução geral é $\sum_i Q_i(n)\lambda_i^n$ sobre as raízes
do \omterm{def:b2:reduction:charpoly}{polinômio característico}) e gasta os dividendos em
aproximação diofantina de $\sqrt2$, na contagem de passeios
e palavras, e num anel de sequências acopladas que só a
diagonalização simultânea consegue desembaraçar.

\begin{problem}[{Problema de fim de semana --- o teorema fundamental das
recorrências lineares}]
\label{pb:b2:reduction:1}
Fixe $k \geq 1$, escalares $a_0, \dots, a_{k-1} \in \C$ com $a_0
\neq 0$, o polinômio mônico $P = X^k - a_{k-1}X^{k-1} - \dots -
a_1 X - a_0$ e a recorrência
\[
(\mathcal R)\colon\quad u_{n+k} = a_{k-1}u_{n+k-1} + \dots +
a_1 u_{n+1} + a_0 u_n \qquad (n \geq 0).
\]
A \emph{matriz companheira} de $P$ é
\[
C =
\begin{pmatrix}
0 & 1 & & \\
 & \ddots & \ddots & \\
 & & 0 & 1\\
a_0 & a_1 & \cdots & a_{k-1}
\end{pmatrix}
\in \mathcal{M}_k(\C).
\]

\medskip
\noindent\textbf{Parte I --- O dicionário companheiro.}
\begin{enumerate}
  \item Mostre que uma sequência $(u_n)$ satisfaz $(\mathcal R)$ se
        e somente se os vetores $v_n = (u_n, u_{n+1}, \dots,
        u_{n+k-1})^{\mathsf T}$ satisfazem $v_{n+1} = Cv_n$, e portanto
        $v_n = C^n v_0$.
  \item Prove que $\chi_C = P$ (desenvolva $\det(XI - C)$ ao longo da
        primeira coluna e faça indução sobre $k$), e depois que também
        $\mu_C = P$ \emph{(passe a $C^{\mathsf T}$, para o qual $e_1$
        é \omterm{def:b2:structures:generated}{cíclico}, e note que uma matriz e sua \omterm{def:b2:linalg:transpose}{transposta} têm
        o mesmo \omterm{def:b2:reduction:polyu}{polinômio minimal})}.
  \item Mostre que, para cada raiz $\lambda$ de $P$, o vetor
        $(1, \lambda, \dots, \lambda^{k-1})^{\mathsf T}$ gera
        o \omterm{def:b2:reduction:eigen}{autoespaço} de $C$ associado a $\lambda$; deduza que
        \emph{todo} \omterm{def:b2:reduction:eigen}{autoespaço} de $C$ tem dimensão $1$ e
        que $C$ é \omterm{def:b2:reduction:diag}{diagonalizável} se e somente se $P$ tem $k$
        raízes distintas.
  \item Suponha que $P$ tenha raízes distintas $\lambda_1, \dots,
        \lambda_k$. Mostre que as sequências geométricas
        $(\lambda_i^n)_n$ formam uma base do espaço das soluções
        de $(\mathcal R)$, de modo que toda solução é $u_n = \sum_i
        c_i\lambda_i^n$ para constantes $c_i$ únicas.
  \item Resolva completamente: $u_{n+2} = u_{n+1} + 6u_n$, $u_0 = 1$,
        $u_1 = 8$.
\end{enumerate}

\medskip
\noindent\textbf{Parte II --- O operador de deslocamento e o
teorema fundamental.} Seja $\mathcal{S}$ o $\C$-espaço vetorial
de todas as sequências complexas e $S \in
\mathcal{L}(\mathcal{S})$ o deslocamento, $S\bigl((u_n)_n\bigr) =
(u_{n+1})_n$.
\begin{enumerate}[resume]
  \item Mostre que o conjunto solução de $(\mathcal R)$ é $\ker
        P(S)$ e que ele tem dimensão exatamente $k$ \emph{(leve uma
        solução a seus valores iniciais)}.
  \item Explique por que o lema da decomposição em núcleos
        (\cref{thm:b2:reduction:kernels}) se aplica a $S$ no
        espaço $\mathcal{S}$ de dimensão infinita sem nenhuma mudança,
        e escreva a decomposição resultante de $\ker P(S)$ para
        $P = \prod_{i=1}^{r}(X - \lambda_i)^{m_i}$ ($\lambda_i$
        distintos, todos não nulos pois $a_0 \neq 0$).
  \item Para $\lambda \neq 0$ e $m \geq 1$, mostre que
        \[
        \ker\,(S - \lambda\,\mathrm{id})^m
        = \bigl\{\,\bigl(Q(n)\,\lambda^n\bigr)_n : Q \in
        \C_{m-1}[X]\,\bigr\},
        \]
        de dimensão $m$. \emph{(Calcule $(S -
        \lambda)\bigl(Q(n)\lambda^n\bigr) =
        \lambda^{n+1}(\Delta Q)(n)$ com $\Delta Q = Q(X + 1) -
        Q(X)$ e use que $\Delta$ baixa o grau; para a
        dimensão, majore-a por $m$ pelos valores iniciais.)}
  \item (O teorema fundamental das recorrências lineares) Conclua:
        se $P = \prod_{i=1}^{r}(X - \lambda_i)^{m_i}$ com os
        $\lambda_i$ distintos e não nulos, as soluções de
        $(\mathcal R)$ são exatamente as sequências
        \[
        u_n = \sum_{i=1}^{r} Q_i(n)\,\lambda_i^n,
        \qquad Q_i \in \C_{m_i - 1}[X],
        \]
        com polinômios $Q_i$ unicamente determinados.
  \item Resolva completamente: $u_{n+2} = 4u_{n+1} - 4u_n$, $u_0 =
        1$, $u_1 = 0$, e confira a resposta em $u_2$.
\end{enumerate}

\medskip
\noindent\textbf{Parte III --- Raízes dominantes e dividendos
diofantinos.}
\begin{enumerate}[resume]
  \item Suponha as raízes simples com $\abs{\lambda_1} >
        \abs{\lambda_i}$ para $i \geq 2$, e $u_n = \sum_i c_i
        \lambda_i^n$ com $c_1 \neq 0$. Mostre que $u_n \sim
        c_1\lambda_1^n$ e $u_{n+1}/u_n \to \lambda_1$.
  \item (Pell) Defina $a_{n+1} = a_n + 2b_n$, $b_{n+1} = a_n +
        b_n$, $a_0 = b_0 = 1$. Mostre que $q(a, b) = a^2 - 2b^2$
        satisfaz $q(a_{n+1}, b_{n+1}) = -q(a_n, b_n)$, logo
        $a_n^2 - 2b_n^2 = (-1)^{n+1}$; relacione isso com o
        \omterm{def:b2:linalg:det}{determinante} de $M = \begin{psmallmatrix}1 & 2\\ 1 &
        1\end{psmallmatrix}$.
  \item Deduza a estimativa de erro
        \[
        \Bigl|\frac{a_n}{b_n} - \sqrt2\Bigr|
        = \frac{1}{b_n\,(a_n + \sqrt2\,b_n)}
        \leq \frac{1}{2b_n^2},
        \]
        e mostre que ela decai geometricamente com razão $3 -
        2\sqrt2$ \emph{(ache os \omterm{def:b2:reduction:eigen}{autovalores} de $M$ e o
        crescimento de $b_n$)}.
  \item (Crescimento geral) A partir da questão 9, prove: (a) se toda
        raiz satisfaz $\abs{\lambda_i} \leq \rho$, então
        $\abs{u_n} \leq C\,n^{m-1}\rho^n$ com $m = \max_i m_i$;
        (b) se existe uma única raiz $\lambda_1$ de módulo
        máximo e $Q_1 \neq 0$, então $u_{n+1}/u_n \to
        \lambda_1$ --- confira isso na solução da questão 10.
\end{enumerate}

\medskip
\noindent\textbf{Parte IV --- Contando passeios e palavras.} Para um
grafo finito com conjunto de vértices $\{1, \dots, N\}$, a
\emph{matriz de adjacência} $A$ tem $A_{ij} = 1$ se $ij$ é uma aresta,
e $0$ caso contrário.
\begin{enumerate}[resume]
  \item Prove que $(A^n)_{ij}$ é o número de passeios de comprimento
        $n$ de $i$ até $j$ (sequências de $n$ arestas, cada passo
        ao longo de uma aresta).
  \item (O triângulo) Para o grafo completo com $3$ vértices,
        $A = J - I$: usando o \omterm{def:b2:reduction:eigen}{espectro} de $J$
        (\cref{ex:b2:linalg:onesmatrix}), mostre que
        \[
        (A^n)_{ii} = \frac{2^n + 2(-1)^n}{3},
        \qquad
        (A^n)_{ij} = \frac{2^n - (-1)^n}{3} \quad (i \neq j),
        \]
        e confira os dois em $n = 2$ listando os passeios.
  \item (Palavras sem $11$) Seja $w_n$ o número de palavras
        binárias de comprimento $n$ sem dois $1$ consecutivos.
        Codifique as palavras por sua última letra para obter uma matriz de
        transferência, mostre que $w_{n+2} = w_{n+1} + w_n$, deduza $w_n =
        F_{n+2}$ (Fibonacci, \cref{exo:b2:reduction:5}) e
        dê a taxa de crescimento $\lim w_{n+1}/w_n$.
  \item (O caminho) Para o grafo caminho $1 - 2 - 3$, mostre que os
        \omterm{def:b2:reduction:eigen}{autovalores} de $A$ são $\sqrt2, 0, -\sqrt2$ com
        \omterm{def:b2:reduction:eigen}{autovetores} $(1, \pm\sqrt2, 1)$ e $(1, 0, -1)$, e
        deduza que o número de passeios de comprimento $n$ de uma ponta
        à outra é $\bigl((\sqrt2)^n + (-\sqrt2)^n\bigr)/4$:
        zero para $n$ ímpar, e $2^{\,n/2 - 1}$ para $n$ par.
        Confira em $n = 4$.
  \item (Fórmula do traço) Mostre que o número total de passeios
        fechados de comprimento $n$ (todos os pontos de partida) é
        $\operatorname{tr}(A^n) = \sum_i \lambda_i^n$, e
        verifique-o no triângulo.
\end{enumerate}

\medskip
\noindent\textbf{Parte V --- Um anel de sequências: diagonalização
simultânea.} Fixe $k \geq 3$, seja $\omega = \eu^{2\iu\pi/k}$
e seja $W \in \mathcal{M}_k(\C)$ o deslocamento \omterm{def:b2:structures:generated}{cíclico}: $W e_i =
e_{i+1}$ (índices módulo $k$, colunas indexadas por $0, \dots, k-1$).
\begin{enumerate}[resume]
  \item Mostre que $W^{\mathsf T}$ é a matriz companheira de
        $X^k - 1$, deduza $\chi_W = \mu_W = X^k - 1$ e que $W$ é
        \omterm{def:b2:reduction:diag}{diagonalizável} com os $k$ \omterm{def:b2:reduction:eigen}{autovalores} simples
        $\omega^j$ e \omterm{def:b2:reduction:eigen}{autovetores} $f_j = (1, \omega^{-j},
        \omega^{-2j}, \dots, \omega^{-(k-1)j})^{\mathsf T}$.
  \item Uma matriz \emph{circulante} é $C = c_0 I + c_1 W + \dots
        + c_{k-1}W^{k-1}$. Mostre que todas as circulantes comutam,
        que a base $(f_0, \dots, f_{k-1})$ diagonaliza
        \emph{todas} elas simultaneamente, e que os
        \omterm{def:b2:reduction:eigen}{autovalores} de $C$ são $\widehat c(\omega^j) = \sum_m
        c_m \omega^{jm}$, $j = 0, \dots, k-1$.
  \item Deduza $\det C = \prod_{j=0}^{k-1} \widehat
        c(\omega^j)$ e verifique que $k = 3$ recupera a
        fatoração de \cref{exo:b2:linalg:8}.
  \item (A média do colar) Seja $x^{(n+1)} = Mx^{(n)}$ com
        $M = \frac12(W + W^{-1})$: cada um de $k$ números dispostos
        em anel é substituído pela média de seus dois
        vizinhos. Mostre que os \omterm{def:b2:reduction:eigen}{autovalores} de $M$ são
        $\cos(2\pi j/k)$ e que o coeficiente de $x^{(0)}$
        em $f_0$ é a média $\frac1k\sum_m x^{(0)}_m$
        \emph{(some as coordenadas dos $f_j$)}.
  \item Conclua: para $k$ ímpar, $x^{(n)}$ converge para o
        vetor constante cujo valor é a média dos valores
        iniciais; para $k = 4$, exiba o \omterm{def:b2:reduction:eigen}{autovalor} responsável
        pela não convergência e a obstrução exata
        (um coeficiente de média alternada que deve se anular).
  \item (Síntese) Em uma frase cada: como a matriz
        companheira converte a análise de $(\mathcal R)$ em
        redução; onde o lema da decomposição em núcleos não precisou
        de dimensão finita; por que os \omterm{def:b2:reduction:eigen}{autovalores} dominantes governam
        as taxas de crescimento e o erro diofantino; por que as potências da
        matriz de adjacência contam passeios; e o que se ganha com
        matrizes que comutam. Nomeie os dois cumes: o teorema
        fundamental das recorrências lineares e --- para as matrizes
        positivas da Parte IV, no volume do terceiro ano de graduação --- o
        teorema de Perron--Frobenius.
\end{enumerate}
\end{problem}
